\section{The annular scalar and its primitive prefix}

For $y>0$ write $\log y_+=\max(\log y,0)$ and define
\[
 R_X=\sum_{n\ge1}\frac{a(n)}{\sqrt n}\log(X/n)_+,
 \qquad
 \A_X=R_X-R_{X/4}.
\]
All sums are finite.  Also put
\[
 H_X(n)=\min\{\log4,\log(X/n)_+\}.
\]
Then directly
\begin{equation}\label{eq:annular-discrete}
 \A_X=\sum_{n\ge1}\frac{a(n)}{\sqrt n}H_X(n).
\end{equation}

\begin{definition}[Primitive prefix]
For $t>0$ set
\[
 C(t)=\sum_{n\le t}\frac{a(n)}{\sqrt n},
 \qquad C(t)=0\quad(0<t<1).
\]
\end{definition}

\begin{theorem}[Exact primitive-prefix window]\label{thm:window}
For every real $X>0$,
\begin{equation}\label{eq:Rwindow}
 R_X=\int_1^X C(t)\frac{dt}{t},
 \qquad
 \boxed{\A_X=\int_{X/4}^{X}C(t)\frac{dt}{t}}.
\end{equation}
\end{theorem}
\begin{proof}
For each $n$,
\[
 \log(X/n)_+=\int_1^X\1_{n\le t}\frac{dt}{t}.
\]
Only $n\le X$ occur, so finite Fubini gives the first identity.  Subtracting
the same identity at $X/4$ gives the second, with the convention $C(t)=0$ for
$t<1$.
\end{proof}

Let
\[
 B_{1/2}(t)=\sum_{n\le t}\frac{\mu(n)}{\sqrt n},
 \qquad B_{1/2}(t)=0\quad(t<1).
\]

\begin{theorem}[Three-band primitive prefix]\label{thm:C-D}
For every real $t>0$,
\begin{equation}\label{eq:Ddef}
 C(t)=6-\D(t),
 \qquad
 \D(t)=6B_{1/2}(t)-\frac9{\sqrt2}B_{1/2}(t/2)
             +\frac32B_{1/2}(t/4).
\end{equation}
\end{theorem}
\begin{proof}
Sum \eqref{eq:aexplicit} over $n\le t$ with weight $n^{-1/2}$.  The three
shifted sums are
\[
 \sum_{2m\le t}\frac{\mu(m)}{\sqrt{2m}}
 =\frac1{\sqrt2}B_{1/2}(t/2),\qquad
 \sum_{4m\le t}\frac{\mu(m)}{\sqrt{4m}}
 =\frac12B_{1/2}(t/4).
\]
The $6\delta_1$ term is $6$ for $t\ge1$ and both sides vanish appropriately
below $1$.
\end{proof}

\begin{corollary}[A stronger pointwise producer]\label{cor:pointwise}
If $C(t)\ge0$ for every sufficiently large $t$, and the finitely many earlier
windows are checked separately, then $\A_X\ge0$ for every sufficiently large
$X$.  In particular, global $C(t)\ge0$ implies global $\A_X\ge0$.
\end{corollary}
\begin{proof}
Use the positive measure $dt/t$ in \eqref{eq:Rwindow}.
\end{proof}

\section{C4MBI67 and exact four-band form}

Substitution of \eqref{eq:aexplicit} into \eqref{eq:annular-discrete} gives
\begin{equation}\label{eq:kappa}
 \A_X=6H_X(1)+\mathfrak C_X,
\end{equation}
where
\[
 \mathfrak C_X=\sum_{n\ge1}\frac{\mu(n)}{\sqrt n}\kappa_X(n),
\qquad
 \kappa_X(n)=-6H_X(n)+\frac9{\sqrt2}H_{X/2}(n)
                         -\frac32H_{X/4}(n).
\]
Thus C4MBI67 is the inequality
\begin{equation}\label{eq:c4mbi}
 \mathfrak C_X\ge-6H_X(1),
\end{equation}
which is exactly $\A_X\ge0$.

\begin{theorem}[Exact four-band boundary]\label{thm:fourband}
For every $X>0$, with $B_{1/2}(t)=0$ below $1$,
\begin{align}
\mathfrak C_X={}&-\frac32\int_{X/16}^{X/8}B_{1/2}(t)\frac{dt}{t}
+\frac{9\sqrt2-3}{2}\int_{X/8}^{X/4}B_{1/2}(t)\frac{dt}{t}\notag\\
&+\frac{9\sqrt2-12}{2}\int_{X/4}^{X/2}B_{1/2}(t)\frac{dt}{t}
-6\int_{X/2}^{X}B_{1/2}(t)\frac{dt}{t}.\label{eq:fourband}
\end{align}
\end{theorem}
\begin{proof}
From Theorems \ref{thm:window} and \ref{thm:C-D},
\[
 \mathfrak C_X=-\int_{X/4}^{X}\D(t)\frac{dt}{t}.
\]
In the second and third terms of \eqref{eq:Ddef}, substitute $u=t/2$ and
$u=t/4$.  Partition the resulting intervals at
$X/8,X/4,X/2$ and collect coefficients.  They are respectively
\[
 -\frac32,\quad -\frac32+\frac9{\sqrt2},\quad
 -6+\frac9{\sqrt2},\quad -6,
\]
which are exactly those in \eqref{eq:fourband}.
\end{proof}

\section{Exact logarithmic prime ownership}

\begin{lemma}[Squarefree logarithmic ownership]\label{lem:owneratom}
For every squarefree integer $n>1$,
\begin{equation}\label{eq:owneratom}
 \mu(n)= -\sum_{p\mid n}\mu(n/p)\frac{\log p}{\log n},
\end{equation}
where the sum is over the distinct prime divisors of $n$.
\end{lemma}
\begin{proof}
Since $n$ is squarefree, $\mu(n/p)=-\mu(n)$ for every $p\mid n$, and
\[
 -\sum_{p\mid n}\mu(n/p)\frac{\log p}{\log n}
 =\mu(n)\frac{\sum_{p\mid n}\log p}{\log n}=\mu(n).
\]
\end{proof}

\begin{theorem}[Prime--Möbius owner decomposition]\label{thm:owner}
For every finitely supported kernel $\kappa$,
\begin{equation}\label{eq:owner}
\sum_{n>1}\frac{\mu(n)}{\sqrt n}\kappa(n)
=-\sum_p\frac{\log p}{\sqrt p}
 \sum_{\substack{m\ge1\\p\nmid m}}
 \frac{\mu(m)}{\sqrt m}\frac{\kappa(pm)}{\log(pm)}.
\end{equation}
\end{theorem}
\begin{proof}
Insert \eqref{eq:owneratom} in the left side, interchange two finite sums, and
write $n=pm$.  The condition $p\nmid m$ is precisely squarefree ownership.
The weights $\log p/\log n$ are nonnegative and sum to one for each source
atom $n$.
\end{proof}

